% Checked exact coface deduction; whole independent audits requested.
\subsection{The two explicit surviving orientation cycles are boundaries}
\label{subsec:k12-survivor-coface}

The orientation survivors above do not obstruct the vanishing of their
homology groups: each has a concrete incoming boundary with coefficient
one.  This is a calculation for the two displayed cells, not a
calculation of the entire Voronoi complexes.

\begin{proposition}[K17-K12-SURVIVOR-BOUNDARY]
\label{prop:k12-survivor-boundary}
Let $\tau$ and $\sigma$ be the cells in
Proposition~\ref{prop:k12-survivor-cycle}, and put
$R=\mathbb Z[1/13!]$. There are oriented Voronoi cells $\eta$ at
rank nine in projective degree twelve and $\eta^+$ at rank ten
in projective degree thirteen for which
\[
             \partial\eta=\tau,\qquad
             \partial\eta^+=\sigma
             \quad\hbox{over }R.                         \tag{CB.1}
\]
The two explicit cycles therefore represent zero in localized
Steinberg homology. No assertion about the remaining classes in
either homology group is made.
\end{proposition}

\begin{proof}
Retain the numbering $v_1,\ldots,v_{12}$ of the minimum-pair
representatives of $\tau$, and set
\[
 w=e_8+e_9,\qquad Q_\eta=H/720,
\]
where
\begingroup
\setlength{\arraycolsep}{2.3pt}
\[
H=\begin{pmatrix}
720&138&140&40&140&0&-180&0&-180\\
138&720&98&0&0&-220&0&-199&80\\
140&98&720&67&-140&0&0&-140&-53\\
40&0&67&720&0&0&-240&0&-93\\
140&0&-140&0&720&-48&-43&-220&-61\\
0&-220&0&0&-48&720&34&-56&-300\\
-180&0&0&-240&-43&34&720&-66&-264\\
0&-199&-140&0&-220&-56&-66&720&-360\\
-180&80&-53&-93&-61&-300&-264&-360&720
\end{pmatrix}.
\]
\endgroup
The standalone exact certificate described below verifies that
$Q_\eta$ is positive definite and that its entire minimum set is
\[
              \{\pm v_1,\ldots,\pm v_{12},\pm w\},
              \qquad \min(Q_\eta)=1.                     \tag{CB.2}
\]
The full minimum set determines a Voronoi cell by the
minimal-class/cell correspondence
\cite[\S1.2.1, p.~4; \S2.1, p.~6]{DEKM2026}.
The thirteen quadratic rays are linearly independent. Hence their
cone is a simplicial Voronoi cone of dimension thirteen and its
projectivization $\eta$ is a degree-twelve cell. Its facets are
exactly the thirteen single-ray deletions. Deleting $w$ gives
$\tau$. Deleting any $v_i$ still gives full vector rank nine.

For each $i$, use the explicit odd integral reflection $g_i$ and
functional $\varphi_i$ in the proof of
Proposition~\ref{prop:k12-survivor-cycle}. Every printed functional
vanishes on $w$. In particular the two expressions involving its
nonzero coordinates are $z_1+z_8-z_9$, which vanish at
$(z_1,z_8,z_9)=(0,1,1)$. Consequently
\[
 g_i(w)=w .
\]
The automorphism $g_i$ therefore stabilizes the facet of $\eta$
obtained by deleting $v_i$, exchanges two of its original
quadratic rays and fixes all its other rays. These rays are
independent, so it reverses the facet orientation. Every one of
these twelve facets is zero in orientation coinvariants over $R$.
The remaining facet $\tau$ has no orientation-reversing stabilizer,
by Proposition~\ref{prop:k12-orientation-survivor}.

We must also check that $\eta$ itself is an allowed orientation
generator. Its stabilizer permutes its facets and preserves the
property of having an orientation-reversing integral stabilizer,
since stabilizers of equivalent cells are conjugate. Thus it
fixes the unique facet without that property, namely $\tau$.
It consequently fixes the complementary ray $\widehat w$ and
restricts on the other twelve rays to a symmetry of $\tau$.
That symmetry has even ray permutation. The permutation on all
thirteen independent rays is therefore even as well. Every
integral stabilizer of $\eta$ preserves its orientation.

For completeness the boundary coefficient is an actual unit,
not an unspecified stabilizer quotient. In the convention of
\cite[\S1.2.1, equations~(1)--(2), p.~4]{DEKM2026}, the coefficient
of an orbit representative is the sum of the oriented incidence
signs over the actual facets in that orbit. There is exactly one
actual facet of $\eta$ which can be equivalent to $\tau$: every
other facet has an odd stabilizer, whereas $\tau$ does not.
The incidence of a facet of a simplex is $\pm1$. With the
orientation on the cone of $\eta$ given by the ordered rays
$(\widehat v_1,\ldots,\widehat v_{12},\widehat w)$ and that on
$\tau$ given by the first twelve rays, the cited convention
makes this coefficient $+1$. Equivalently, the equivariant
cellular boundary has transfer index
$[\Gamma_\eta:\Gamma_\eta\cap\Gamma_\tau]=1$, because the facet
is fixed by $\Gamma_\eta$; the subsequent inclusion into
$\Gamma_\tau$ contributes no multiplier on degree-zero
coinvariants. There is no additional target-stabilizer factor.
This proves the first equality in (CB.1).

For the second, take the orthogonal extension
$Q_\eta\perp[1]$ on $\mathbb Z^{10}$. Its minimum pairs are
the thirteen in (CB.2), embedded in the first nine coordinates,
and $\pm e_{10}$: a vector with both components nonzero has
norm at least two. Its fourteen quadratic rays are independent,
and the projective degree is thirteen. Call this cell $\eta^+$.
Deleting $\widehat w$ gives $\sigma$. Deleting any original
$\widehat v_i$ gives a full-rank facet whose odd stabilizer is
$g_i\perp1$. Deleting $\widehat e_{10}$ leaves vector rank nine
inside the rank-ten ambient lattice, so that facet belongs to
the relative boundary subcomplex and contributes zero.

The unique full-rank facet of $\eta^+$ with a nonzero orientation
generator is therefore $\sigma$. Every stabilizer of $\eta^+$
preserves it and fixes its complementary ray, so the preceding
argument also proves that $\eta^+$ is orientation preserving.
The same incidence calculation gives coefficient $\pm1$ on
$\sigma$ and no other nonzero term. Choose the orientation of
$\eta^+$ to make that coefficient $+1$. This proves (CB.1).
\end{proof}

\paragraph{Complete exact verification and its limits.}
The self-contained program
\texttt{k12\_survivor\_coface\_certificate.py} contains the displayed
integer matrix, the thirteen vectors and all twelve reflection
rows as literal integer data. It uses only rational arithmetic,
integer square roots and the Python standard library.
It factors $Q_\eta=L D L^t$ exactly, checks every diagonal entry
of $D$ is positive and checks the matrix identity itself.
For enumeration it uses
\[
 Q_\eta[z]=\sum_i D_i
       \left(z_i+\sum_{j>i}L_{ji}z_j\right)^2.
\]
Descending in $i$, if the remaining budget is $B$ and the shift
is $a/b$ with $b>0$, the next coordinate is restricted by
\[
 (b z_i+a)^2\le (B/D_i)b^2.
\]
An integer square root therefore gives the exact complete
interval of possible integers $z_i$; every such integer is
visited. The executed search visits $388$ nodes and finds
exactly the twenty-six signed vectors in (CB.2), with no
other nonzero integer vector of norm at most one.
The program checks their norms directly, the rank thirteen of
their quadratic rays, and the full vector and quadratic ranks
of every other facet. It then checks the twelve integral
involutions, their transpositions of facet rays, and the
identities $\varphi_i(w)=0$. The complete rational pivots,
minimum list and witness matrices are retained in the
accompanying JSON output.

Floating-point linear programming was used only to discover a
candidate; it is not an input to this final verifier. The boundary
argument uses the formal Voronoi chain convention, not the
questioned endpoint enumeration in Theorem~4.5 of the same
source. These calculations kill the two specified cycles.
They neither establish the whole rank-nine or rank-ten target
homology vanishing nor supply the missing rank-eight
$H_5(\mathrm{GL}_8;\mathrm{St}_8\otimes R)$ input to the
$K_{12}$ criterion.
